\section*{Chapitre \ref{ch:g10:earth-energy-budget} --- Le bilan énergétique de la Terre}

\begin{solution}{exo:g10:earth-energy-budget:1}
L'\omterm{def:g10:earth-energy-budget:insolation}{insolation} d'un lieu est la puissance de la lumière solaire reçue par un
mètre carré horizontal au sommet de l'atmosphère, au-dessus de ce lieu
(moyennée sur une journée ou sur une année). L'\omterm{def:g10:earth-energy-budget:albedo}{albédo} d'une surface est la
fraction de la lumière solaire incidente qu'elle réfléchit. La neige fraîche
(0.8 à 0.9) réfléchit bien plus que l'océan libre (0.06).
\end{solution}

\begin{solution}{exo:g10:earth-energy-budget:2}
Réfléchie~: $100/340 \approx \qty{29}{\percent}$~; absorbée par l'air~:
$80/340 \approx \qty{24}{\percent}$~; absorbée par la surface~: $160/340
\approx \qty{47}{\percent}$.
\end{solution}

\begin{solution}{exo:g10:earth-energy-budget:3}
Asphalte~: $800 \times (1 - 0.04) = \qty{768}{W}$. Neige~: $800 \times (1 -
0.85) = \qty{120}{W}$, plus de six fois moins.
\end{solution}

\begin{solution}{exo:g10:earth-energy-budget:4}
L'évaporation~: \qty{82}{W/m^2}, contre \qty{21}{W/m^2} pour le
réchauffement de l'air. L'eau qui s'évapore prend sa \omterm{def:g10:earth-energy-budget:latent-heat}{chaleur latente} à la
mer~; la vapeur monte avec l'air, et là où elle se condense en gouttelettes
de nuage, souvent plusieurs kilomètres plus haut, elle rend cette énergie à
l'air qui l'entoure.
\end{solution}

\begin{solution}{exo:g10:earth-energy-budget:5}
En décembre, l'hémisphère Nord penche à l'opposé du Soleil~: au pôle, le
Soleil reste sous l'horizon toute la journée, donc aucune lumière solaire
n'arrive. Au solstice de décembre, c'est le cas partout au-delà de
$90^\circ - 23.44^\circ =
\qty{66.56}{\degree}$ N.
\end{solution}

\begin{solution}{exo:g10:earth-energy-budget:6}
$1361 \sin 90^\circ = \qty{1361}{W}$~; $1361 \sin 30^\circ =
\qty{680.5}{W}$~; $1361 \sin 10^\circ \approx \qty{236}{W}$. La moitié de
$S_0$ demande $\sin h = 0.5$, soit $h = \qty{30}{\degree}$.
\end{solution}

\begin{solution}{exo:g10:earth-energy-budget:7}
Réfléchis 100, absorbés par l'air 80, absorbés par la surface 160~:
$100 + 80 + 160 = 340$, toute la lumière solaire est prise en compte. À la
surface, $185 = 25 + 160$~: ce qui n'est pas réfléchi est absorbé.
\end{solution}

\begin{solution}{exo:g10:earth-energy-budget:8}
Disque~: $\pi \times (6.371\times10^6)^2 \approx \qty{1.275e14}{m^2}$~;
puissance~: $1361 \times 1.275\times10^{14} \approx \qty{1.74e17}{W}$.
Sphère~: $4\pi
R^2 \approx \qty{5.10e14}{m^2}$. Moyenne~: $1.74\times10^{17}/5.10\times
10^{14} \approx \qty{340}{W/m^2}$, le quart de $S_0$.
\end{solution}

\begin{solution}{exo:g10:earth-energy-budget:9}
$0.7 \times 5.10\times10^{14} \times 3.16\times10^7 \approx
\qty{1.1e22}{J}$ en une année.
\end{solution}

\begin{solution}{exo:g10:earth-energy-budget:10}
Les nuages réfléchissent $100 - 53 = \qty{47}{W/m^2}$ de lumière solaire en
plus, qui n'est donc pas absorbée. Ils retiennent $267 - 239 = \qty{28}{W/m^2}$ d'infrarouge. Au total, la Terre perd
$47 - 28 = \qty{19}{W/m^2}$~: les nuages, dans leur ensemble, la
refroidissent.
\end{solution}

\begin{solution}{exo:g10:earth-energy-budget:11}
D'environ le 10~mai à la fin juillet, c'est-à-dire en mai, juin et juillet.
Sur l'équateur, le Soleil passe au zénith à midi aux deux équinoxes, en mars
et en septembre~; aux solstices, il passe à \qty{23.44}{\degree} du zénith,
si bien que l'équateur reçoit un peu moins en juin et en décembre.
\end{solution}

\begin{solution}{exo:g10:earth-energy-budget:12}
Chaque seconde~: $82/(2.47\times10^6) \approx \qty{3.32e-5}{kg}$. Chaque
jour~: $3.32\times10^{-5} \times 86\,400 \approx \qty{2.9}{kg}$, une couche
d'environ \qty{2.9}{mm}. Chaque année~: $3.32\times10^{-5} \times 3.16\times10^7 \approx
\qty{1050}{kg}$, environ \qty{1.05}{m} d'eau.
\end{solution}

\begin{solution}{exo:g10:earth-energy-budget:13}
$S_0 \sin\delta = S_0/\pi$ donne $\sin\delta = 1/\pi \approx 0.318$, donc
$\delta \approx \qty{18.6}{\degree}$. Alors $\sin\lambda = 0.318/\sin 23.44^\circ
\approx 0.318/0.398 \approx 0.800$~: $\lambda \approx \qty{53.15}{\degree}$ ou
$180^\circ - 53.15^\circ = \qty{126.85}{\degree}$. Entre les deux, $\lambda$
croît de \qty{73.7}{\degree}, soit $73.7/360$ d'année, environ \num{75}
jours, de la mi-mai à la fin juillet.
\end{solution}

\begin{solution}{exo:g10:earth-energy-budget:14}
La glace absorbe $1 - 0.6 = 0.40$ de la lumière solaire, l'océan $1 - 0.06 =
0.94$. Puissance supplémentaire~: $(0.94 - 0.40) \times 300 \times 10^{12} \approx
\qty{1.6e14}{W}$, plus que toute la photosynthèse de la planète. L'énergie
supplémentaire réchauffe l'eau, qui fait fondre davantage de glace au bord
des plaques, ce qui découvre davantage d'eau sombre~: le changement
s'entretient lui-même.
\end{solution}

\begin{solution}{exo:g10:earth-energy-budget:15}
Au plus près~: $1361 \times (149.6/147.1)^2 \approx \qty{1408}{W/m^2}$~; au
plus loin~: $1361 \times (149.6/152.1)^2 \approx \qty{1317}{W/m^2}$~; un
écart d'environ \qty{7}{\percent}. À \qty{60}{\degree} N, l'\omterm{def:g10:earth-energy-budget:insolation}{insolation} est
vingt fois plus forte en juin qu'en décembre~: une variation de quelques
pour cent ne peut pas en être la cause. De plus, la Terre est au plus près
du Soleil début janvier, pendant l'hiver boréal, et les deux hémisphères, à
la même distance du Soleil, ont des saisons opposées~: c'est l'inclinaison
de l'axe, et non la distance, qui fait les saisons.
\end{solution}

\begin{solution}{pb:g10:earth-energy-budget:1}
\textbf{1.} $\pi \times (6.371\times10^6)^2 \approx \qty{1.275e14}{m^2}$.

\textbf{2.} $1361 \times 1.275\times10^{14} \approx \qty{1.735e17}{W}$.

\textbf{3.} $4\pi R^2 \approx \qty{5.101e14}{m^2}$.

\textbf{4.} $1.735\times10^{17}/5.101\times10^{14} \approx \qty{340}{W/m^2}$,
soit $S_0/4$.

\textbf{5.} Équateur~: \qty{1361}{W}~; \qty{60}{\degree} N~: $1361 \sin
30^\circ = \qty{680.5}{W}$.

\textbf{6.} $100/340 \approx 0.294$.

\textbf{7.} $340 - 100 = \qty{240}{W/m^2}$~; sur toute la Terre, $240 \times
5.101\times10^{14} \approx \qty{1.224e17}{W}$.

\textbf{8.} $53/340 \approx 0.156$. Les nuages multiplient l'\omterm{def:g10:earth-energy-budget:albedo}{albédo} par
$100/53 \approx 1.9$~: ils le doublent presque.

\textbf{9.} $(240/5.67\times10^{-8})^{1/4} \approx \qty{255}{K}$, environ
\qty{33}{K} de moins que les \qty{288}{K} de la surface.

\textbf{10.} $160 + 342 = \qty{502}{W/m^2}$.

\textbf{11.} $82 + 21 + 398 = \qty{501}{W/m^2}$. La surface gagne environ
\qty{1}{W/m^2}~: le déséquilibre de \qty{0.7}{W/m^2}, arrondi.

\textbf{12.} $342/160 \approx 2.1$. L'air envoie à la surface, sous forme
d'infrarouge, plus de deux fois l'énergie de la lumière solaire qu'elle
absorbe~: c'est l'effet de serre.

\textbf{13.} $82 \times 3.156\times10^7/(2.47\times10^6) \approx
\qty{1050}{kg}$, une couche d'environ \qty{1.05}{m} d'épaisseur.

\textbf{14.} $0.7 \times 5.101\times10^{14} \approx \qty{3.6e14}{W}$.

\textbf{15.} $2805/72.07 \approx \qty{38.9}{kJ}$ par gramme de carbone.

\textbf{16.} $104.9\times10^9$ tonnes font $1.049\times10^{17}$~g~;
$1.049\times10^{17} \times 38.9\times10^3 \approx \qty{4.08e21}{J}$.

\textbf{17.} $4.08\times10^{21}/3.156\times10^7 \approx \qty{1.29e14}{W}$.

\textbf{18.} $1.29\times10^{14}/5.101\times10^{14} \approx
\qty{0.25}{W/m^2}$.

\textbf{19.} $0.7/0.25 \approx 2.8$~: le déséquilibre de la Terre vaut
presque trois fois la puissance que stocke la photosynthèse.

\textbf{20.} $1.29\times10^{14}/1.224\times10^{17} \approx 0.0011$~: la
photosynthèse stocke \textbf{environ \qty{0.1}{\percent}} de la lumière
solaire absorbée par la Terre --- un millième.
\end{solution}
